\chapter{Validation against ergm}\label{ch:validation}

\nwergm{} is a reimplementation, and the only reason to trust its numbers is that they agree with the program that it reimplements. This chapter describes how the comparison with \ergm{}~\ergmver{} was done, shows a comparison that was run for this manual, and lists the quirks of \ergm{} that \nwergm{} reproduces and the places where the two differ. The status of every term is in Chapter~\ref{ch:terms}.

\section{How the comparison was done}\label{sec:howvalidated}

The reference is \ergm{}~\ergmver{} (R~4.6). The comparisons are of four kinds, from the strictest to the loosest. The first three are scripted and rerun by the test suite of the package, 60 do-files in the directory \texttt{cscripts}; the generators of the R reference values are in the directory \texttt{dev}.

\begin{enumerate}[nosep]
  \item \textbf{Observed statistics.} For every term, the observed statistic of test networks is computed by \nwergm{} and by \texttt{summary()} of \ergm{}. The test networks are random undirected, directed and two-mode networks with 14 to 100 nodes and the real networks of the statnet examples (the Florentine marriages, Sampson's monastery, the Kapferer tailor shop, Mesa High School). The statistics agree exactly, or to the floating-point precision of the statistic (at most $2\times10^{-7}$, for the term that \ergm{} computes in single precision).
  \item \textbf{Change statistics and pseudolikelihood.} The change statistics determine the pseudolikelihood. The change statistic of every dyad is compared with the difference of the statistic computed by brute force with and without the tie, and, for the two-mode terms, with the design matrix of \texttt{ergmMPLE()}. The pseudolikelihood coefficients and standard errors are compared with \texttt{ergm(estimate="MPLE")} on the same networks: the coefficients agree to $10^{-4}$ or better for every term that was compared, and to $10^{-6}$ to $10^{-10}$ for most (the column \emph{Validation} of Chapter~\ref{ch:terms} gives the figure per term), with the exceptions that Sections~\ref{sec:quirks} and~\ref{sec:differs} name. Where R's own fit is degenerate or separated (an infinite or missing coefficient, a standard error above 20), the comparison is not possible and the case is skipped.
  \item \textbf{Monte Carlo maximum likelihood.} For dyad-dependent models, the fit by \nwergm{} is repeated with several seeds and compared with several runs of \ergm{} on the same network. The criterion is that the difference of the mean coefficients is small compared with the standard error and with the run-to-run spread of both programs: for the models of the structural terms every coefficient was within 0.1 posterior standard errors of R's and inside the spread over seeds; for two-mode terms the differences of the means were between 0.01 and 0.14 standard errors. Constrained fits and fits with missing dyads were compared in the same way.
  \item \textbf{The sampler itself.} The Markov chains were checked against exact enumeration of the distribution on networks that are small enough (5-node undirected, 4-node directed and $3\times3$ two-mode networks; with every constraint that restricts the sample space) and against \texttt{simulate()} of \ergm{} (means and standard deviations of the statistics of long chains agree, see below).
\end{enumerate}
The estimation controls (the Hummel step length, the convergence test, the standard errors) were ported from the algorithms of \ergm{} and checked against intermediate results of R: the step-length routine against \ergm{}'s function to $10^{-7}$, the multivariate spectral variance estimate used by \texttt{mcmc.diagnostics()} to $10^{-8}$, the effective sample size and the Geweke and Heidelberger--Welch statistics against \texttt{coda}.

\section{A comparison run for this manual}\label{sec:cmpmanual}

To show what these comparisons look like, the main ones were repeated for this manual on the networks of Chapter~\ref{ch:examples}. The Stata side is the do-file \texttt{manual/ex/ex91\_compare\_stata.do}, the R side \texttt{manual/ex/r/compare.R}; both read the same networks and attributes from files.

\paragraph{Observed statistics.} Table~\ref{tab:cmpstat} compares the observed statistics of the Lazega network. The statistics of the counts agree exactly; the geometrically weighted statistic, which is a sum of non-integer weights, agrees to \cmpStatMax.

\paragraph{Pseudolikelihood.} Table~\ref{tab:cmpmple} compares the pseudolikelihood estimates of four models: a dyad-independent model of the Florentine network, a dyad-dependent model (with \texttt{mutual}) of the Lazega network, and two models of the simulated two-mode network (one dyad-independent with \texttt{nodemix}, one dyad-dependent with two-mode degree and star terms). The largest difference between the coefficients is \cmpMpleMax{}; the standard errors agree to a relative difference of at most \cmpSeRel.

\paragraph{Monte Carlo maximum likelihood.} Table~\ref{tab:cmpmcmle} compares five runs of each program on the Lazega model of Section~\ref{sec:lazega}. With its default settings, which since version 1.2.0 are those of \ergm{} (burn-in 16384, interval 1024), \nwergm{} has a spread over seeds of the order of R's (standard deviations of the coefficients over the five seeds between 0.000 and 0.006 for both programs), and the means differ from R's by at most \cmpMcmleSd{} of a standard error; the standard errors of the two programs agree. The middle columns repeat the fits with the settings of the versions before 1.2.0 (burn-in 3000, interval 50): the chains are then short relative to the autocorrelation (Section~\ref{sec:lazega}), the spread over seeds is up to ten times larger (0.033 for \opt{mutual}, against 0.005 with the new defaults and 0.006 for \ergm{}), and single fits are off by a larger part of a standard error.

\paragraph{The sampler.} Table~\ref{tab:cmpsim} compares the means and standard deviations of the five model statistics in chains simulated at the coefficients of the Lazega fit, by \nwergm{} and by \texttt{simulate()} of \ergm{}. The means differ by at most \cmpSimRel{} of a standard deviation and the standard deviations agree to the first digits.

\begin{table}[ht]\centering\small\caption{Observed statistics of the Lazega friendship network: nwergm (the first row of the model block of \texttt{estat gof}) and \texttt{summary()} of ergm.}\label{tab:cmpstat}
\begin{tabular}{@{}lrrr@{}}\toprule
term & nwergm & ergm & difference\\\midrule
\opt{edges} & 575 & 575 & 0.0e+00\\
\opt{mutual} & 74 & 74 & 0.0e+00\\
\opt{nodematch\_prac} & 412 & 412 & 0.0e+00\\
\opt{nodeicov\_seniority} & 19379 & 19379 & 0.0e+00\\
\opt{gwesp\_0.5} & 555.6191 & 555.6191 & 1.5e-08\\
\bottomrule\end{tabular}\end{table}
\begin{table}[ht]\centering\small\caption{Pseudolikelihood estimates: nwergm (\texttt{estimate(mple)}) and ergm (\texttt{estimate="MPLE"}) on the same data.}\label{tab:cmpmple}
\begin{tabular}{@{}llrrrrr@{}}\toprule
fit and term & & nwergm & ergm & $|\Delta|$ & SE nwergm & SE ergm\\\midrule
\multicolumn{7}{@{}l}{\emph{Florentine marriages}}\\
& \opt{edges} & -2.59493 & -2.59493 & 2.1e-10 & 0.53606 & 0.53606\\
& \opt{nodecov\_wealth} & 0.01055 & 0.01055 & 3.8e-12 & 0.00467 & 0.00467\\
\multicolumn{7}{@{}l}{\emph{Lazega friendship}}\\
& \opt{edges} & -2.50763 & -2.50763 & 4.1e-08 & 0.11080 & 0.11079\\
& \opt{mutual} & 0.99645 & 0.99645 & 2.2e-08 & 0.11050 & 0.11050\\
& \opt{nodematch\_prac} & 0.91390 & 0.91390 & 1.1e-08 & 0.09903 & 0.09902\\
& \opt{nodeicov\_seniority} & -0.00691 & -0.00691 & 2.6e-10 & 0.00224 & 0.00224\\
\multicolumn{7}{@{}l}{\emph{simulated two-mode network, model 1}}\\
& \opt{edges} & -0.75633 & -0.75633 & 1.9e-11 & 0.17873 & 0.17873\\
& \opt{nodemix\_kind\_2\_1} & -1.51236 & -1.51236 & 2.8e-09 & 0.39313 & 0.39312\\
& \opt{nodemix\_kind\_1\_2} & -2.37917 & -2.37917 & 2.4e-08 & 0.61625 & 0.61620\\
& \opt{nodemix\_kind\_2\_2} & -0.23407 & -0.23407 & 4.6e-11 & 0.37073 & 0.37073\\
\multicolumn{7}{@{}l}{\emph{simulated two-mode network, model 2}}\\
& \opt{edges} & -0.25048 & -0.25048 & 4.1e-11 & 0.37254 & 0.37254\\
& \opt{nodemix\_kind\_2\_1} & -1.60452 & -1.60452 & 9.9e-11 & 0.42091 & 0.42091\\
& \opt{nodemix\_kind\_1\_2} & -4.02396 & -4.02396 & 1.1e-09 & 0.76838 & 0.76837\\
& \opt{nodemix\_kind\_2\_2} & -1.68223 & -1.68223 & 7.4e-10 & 0.52610 & 0.52610\\
& \opt{b2degree\_5} & 2.64228 & 2.64228 & 2.4e-10 & 0.51729 & 0.51728\\
& \opt{b2degree\_6} & 2.75054 & 2.75054 & 6.1e-10 & 0.62388 & 0.62387\\
& \opt{b1star\_2} & -0.13954 & -0.13954 & 1.7e-11 & 0.13553 & 0.13553\\
\bottomrule\end{tabular}\end{table}
\begin{table}[ht]\centering\small\caption{Monte Carlo maximum likelihood, Lazega friendship network, five seeds each: mean and (in parentheses) standard deviation over the seeds of the coefficients, and mean standard error. nwergm with its default settings (adaptive sample), nwergm with the settings of version 1.1 and fixed-size samples (\texttt{mcmleeffectivesize(0) mcmcburnin(3000) mcmcinterval(50)}: 3000 draws per iteration), and ergm with its default settings.}\label{tab:cmpmcmle}
\begin{tabular}{@{}lrrrrrr@{}}\toprule
 & \multicolumn{2}{c}{nwergm, default} & \multicolumn{2}{c}{nwergm, fixed, old settings} & \multicolumn{2}{c}{ergm}\\\cmidrule(lr){2-3}\cmidrule(lr){4-5}\cmidrule(lr){6-7}
term & mean (sd) & SE & mean (sd) & SE & mean (sd) & SE\\\midrule
\opt{edges} & -3.318 (0.010) & 0.133 & -3.332 (0.009) & 0.138 & -3.315 (0.003) & 0.136\\
\opt{mutual} & 0.804 (0.008) & 0.160 & 0.817 (0.033) & 0.174 & 0.804 (0.006) & 0.161\\
\opt{nodematch\_prac} & 0.714 (0.007) & 0.081 & 0.714 (0.029) & 0.087 & 0.715 (0.004) & 0.081\\
\opt{nodeicov\_seniority} & -0.004 (0.000) & 0.002 & -0.004 (0.001) & 0.002 & -0.005 (0.000) & 0.002\\
\opt{gwesp\_0.5} & 0.488 (0.004) & 0.064 & 0.495 (0.009) & 0.064 & 0.488 (0.002) & 0.065\\
\bottomrule\end{tabular}\end{table}
\begin{table}[ht]\centering\small\caption{Simulation at given coefficients (the coefficients of the Lazega fit): mean and standard deviation of the model statistics over three chains of 4000 draws, nwergm (\texttt{nwergm simulate} at the fitted coefficients with \texttt{burnin(20000) interval(1000)}) and ergm (\texttt{simulate()} with its default burn-in 16384 and interval 1024).}\label{tab:cmpsim}
\begin{tabular}{@{}lrrrr@{}}\toprule
term & mean nwergm & mean ergm & sd nwergm & sd ergm\\\midrule
\opt{edges} & 578.6 & 579.9 & 40.3 & 39.5\\
\opt{mutual} & 74.8 & 75.1 & 10.5 & 10.5\\
\opt{nodematch\_prac} & 416.1 & 416.8 & 31.2 & 30.2\\
\opt{nodeicov\_seniority} & 19464.5 & 19521.3 & 1497.4 & 1466.7\\
\opt{gwesp\_0.5} & 564.6 & 566.9 & 76.1 & 74.8\\
\bottomrule\end{tabular}\end{table}


\section{Quirks of ergm that nwergm reproduces}\label{sec:quirks}

Agreeing with \ergm{} sometimes means agreeing with its idiosyncrasies. \nwergm{} reproduces the following, which are documented here so that a difference from the theory in the literature is not mistaken for a mistake in \nwergm{}.
\begin{itemize}
  \item \textbf{\texttt{triangle} on a directed network.} In \ergm{}'s change statistic a node that is both an in- and an out-neighbor of the head is visited twice and the change adds the ties in both directions; \nwergm{} copies this (statistic and pseudolikelihood agree with R).
  \item \textbf{Single precision.} \texttt{tripercent} computes its ratios in single precision in \ergm{}, and its observed statistic is the sum of the single-step changes along a build-up of the network; \nwergm{} copies both, and agrees with R's statistic to $2\times10^{-7}$.
  \item \textbf{Scaling by the observed network.} \texttt{degcrossprod} and \texttt{degcor} are scaled by quantities of the observed network at the moment of model creation (twice the number of ties, the variance of the observed degrees), and for \texttt{degcor} the summary statistic is the true degree correlation whereas the change statistic is an approximation. \nwergm{} copies this; the terms are therefore not available in \texttt{nwergm simulate}, where there is no observed network.
  \item \textbf{\texttt{Prod()}.} \ergm{} computes the product as \texttt{Exp(Sum(Log()))} with $\log0=-1/\sqrt{\epsilon}$. Built up tie by tie from the empty network, the statistic passes through values of about $7\times10^{7}$ and loses eight digits: for Mesa High School R gives 522.48718785 for a product whose exact value is $(3491\cdot659\cdot62)^{1/3}=522.48719458$, the value that \nwergm{} computes. The tests of \texttt{Prod()} therefore use the tolerance $10^{-6}$ for the statistics; the pseudolikelihood coefficients agree to $10^{-5}$.
  \item \textbf{\texttt{Symmetrize(rule="weak")}.} \texttt{summary()} of a formula with this operator is right, because \ergm{} adds the arcs one by one. But \texttt{ergm()} and \texttt{ergmMPLE()} start from the observed network and fill the auxiliary network with an initialization rule that, for \texttt{weak}, toggles an arc unless its reverse exists, so that every reciprocated dyad is missing from the symmetrized network on which the pseudolikelihood change statistics are computed (108 of 182 change statistics differed from brute-force differences of \texttt{summary()} on a 14-node directed network; the rules \texttt{strong}, \texttt{upper} and \texttt{lower} agree). \nwergm{} computes the exact change statistics, so that its pseudolikelihood equals R's only on a network without reciprocated ties. The test therefore compares \texttt{rule(weak)} on networks with reciprocated ties with a pseudolikelihood fitted to R's own \texttt{summary()} differences, and on a network without reciprocated ties with \texttt{ergm(estimate="MPLE")} itself. This is not to be ``fixed'' to reproduce R's initial network.
  \item \textbf{\texttt{opentriad}.} The change statistic of \ergm{}~\ergmver{} overflows an unsigned integer when a tie is removed, so that its estimation with this term is wrong (the pseudolikelihood coefficient comes out as 0 with a standard error of $10^{-7}$) although its \texttt{summary()} is right. \nwergm{} uses the correct change statistic and therefore differs from R in the estimates for this term.
  \item \textbf{\texttt{gwesp} and \texttt{esp} with \texttt{type="RTP"}.} \ergm{}~\ergmver{} returns wrong values for the reciprocated two-path definition by default (a defect that is corrected in the development version); \nwergm{} agrees with R when R is run with \texttt{term.options=list(cache.sp=FALSE)}, to $10^{-7}$.
  \item \textbf{The penalized pseudolikelihood.} In \ergm{}~\ergmver{} \texttt{MPLE.type="penalized"} stops with an error under R~4.3 or later (a logical of length 3 in a condition); the development version fixes it. \nwergm{} reproduces the fixed behavior, which was compared with the development version to $10^{-6}$ on five networks, one of them separated.
  \item \textbf{\texttt{mcmc.diagnostics()}.} The multivariate spectral variance used for the Hotelling omnibus test and the joint Geweke test is that of \ergm{}~\ergmver{} (a Yule--Walker autoregression); the development version uses an ordinary-least-squares vector autoregression, which changes the second or third digit of these $p$-values.
  \item \textbf{Constraints.} R refuses \texttt{blockdiag} on non-contiguous blocks under the default \texttt{noncontig="merge"}, and \nwergm{} too. The constraint \texttt{hamming} is broken in R and refused; \texttt{odegreedist} on a directed network and the degree-distribution constraint on two-mode networks fail in R's released code and are refused. R's \texttt{degrees} and \texttt{degreedist} proposals for directed networks are not exactly symmetric; \nwergm{} reproduces the same chain, which agrees with R's \texttt{simulate()} within Monte Carlo error but not exactly with an exact enumeration (by at most 1\% on a 4-node directed network).
  \item \textbf{Contrastive divergence.} \ergm{} reports the covariance of a contrastive-divergence estimate as missing; \nwergm{} stores the covariance of the sample (Section~\ref{sec:cd}).
  \item \textbf{Separation.} Where a level has no node or is perfectly predicted, R reports \texttt{-Inf} or \texttt{NA} and \nwergm{} reports 0 with a standard error of 0 (for example for \texttt{idegree} and \texttt{sender}).
\end{itemize}

\section{Where nwergm differs from ergm}\label{sec:differs}

The differences that matter to a user are these.
\begin{itemize}
  \item \textbf{Monte Carlo defaults.} Burn-in 16384 and interval 1024 are R's defaults (before version 1.2.0: 3000 and 50). \nwergm{} records a fixed 3000 draws, where \ergm{} adapts the number of draws to an effective sample size of 64 (at least 256 draws); both stop after at most 60 iterations, and \nwergm{} lets the interval and the sample size grow when the effective sample size is below the larger of 200 and 64 times the number of coefficients (a stricter floor than R's 64). The pseudo-random streams differ.
  \item \textbf{Curved fits on boundary solutions.} Where the data do not identify the decay of a curved term, both programs end at a boundary, and the end points can differ: in a comparison with random networks, the fits of \texttt{gwespfree} (directed, OTP), \texttt{gwidegreefree}, \texttt{gwdspfree} and \texttt{gwnspfree} ended at a decay of about $10^{-12}$ in R and at another small value, $10^{-6}$ or less, in \nwergm{}. Where the data identify the curve the fits agree to $10^{-8}$ to $10^{-6}$.
  \item \textbf{\texttt{gwdegree} on a directed network.} R refuses it; \nwergm{} accepts it, with no counterpart in R to check the result against.
  \item \textbf{Controls not available.} \texttt{MCMLE.metric} (the log-normal approximation is the only one), \texttt{MCMLE.dampening}, \texttt{MCMLE.sequential}, \texttt{MCMLE.density.guard}, \texttt{parallel}, \texttt{monitor=}, \texttt{response=}, \texttt{esteq=}, and the constraint \texttt{constraints=}-arguments \texttt{ChangeStats()}, \texttt{Dyads()} and \texttt{I()}.
  \item \textbf{Proposals not available.} The Hamming proposals (reachable only through the broken \texttt{hamming} constraint), the shared-partner proposal for the \texttt{triadic} hint, \texttt{TNT10} (which R marks as flawed), the bipartite degree-distribution proposal (whose code writes the second toggle into the wrong slot), the degree-mixing proposal, and the stratified proposal of the \texttt{strat} hint. The bounded-degree proposal is realized as the tie/no-tie proposal with a rejection of moves that violate the bounds, which has the same target distribution.
  \item \textbf{Unbounded loops.} The degree-preserving proposals of R can loop forever on a state without a valid tetrad (for example on 4-node directed networks); \nwergm{} bounds the search at $10^{6}$ tries and treats the step as a null move.
  \item \textbf{Terms.} The arguments listed in the remarks of Chapter~\ref{ch:terms} that are not available, in particular \texttt{levels=} for the terms of \opt{structural()}; the form of \texttt{S()} with two sets of nodes; and the \texttt{Curve()} operator beyond linear maps.
  \item \textbf{Missing dyads and constraints together.} Missing dyads cannot be combined with \opt{freedyads()}, \opt{blockdiag()} or \opt{fixdensity}; an observation constraint in general is available through \opt{obsconstraints()}.
\end{itemize}
